%%%%%%%%%%%%%%%%%%%%%%%%%%%%%%%%%%%%%%%%%%%%%%%%%%%%
\section{はじめに}

自律AIシステムおよびロボットエージェントは、行動パターンを内包したセンサイベントストリームを生成する。
これらのパターンが「いつ」密集して共起するかを理解すること、単にグローバルに頻出するかどうかではなく、
は、知的環境におけるリアルタイム監視・異常検知・適応制御に不可欠である。
マルチエージェントシステムやセンサ搭載ロボットプラットフォームでは、特定のイベント集合が特定の時間区間に
密集していることを検出することで、協調行動・日常ルーティン・迅速な対応を要する異常の出現を
明らかにすることができる。
また、この手法は小売や疫学などの幅広い分野へも適用可能である。

本稿では、スキャン頻度 $\Delta$ をパラメータとする密集区間マイニングアルゴリズム
\textit{Streaming Apriori-Window} を提案する。
アルゴリズムは、固定幅スライディングウィンドウ内でのアイテムセット出現頻度が閾値を満たす
すべての時間区間を特定する。
2つの高速化技術、Scan Stride Skipping と Apriori Interval Pruning を組み込み、
実用的な処理速度を実現する。
著者らは予備的研究~\cite{takayasu2025arxiv}においてApriori-Windowを提案した。
本稿では、これをストリーミングフレームワークへ拡張し、形式的な出力保証、高速化技術の証明、
および包括的な実験評価を加える。

本稿の主な貢献は以下の通りである：
\begin{itemize}
  \item 出力レイテンシと処理コストのトレードオフを制御する明示的なスキャン頻度パラメータ $\Delta$ を持つ
        ストリーミング密集区間マイニング問題を形式化する。
  \item 確定出力保証を示す：現在のスキャン位置より前に終端する区間は、将来のイベントによって変化しない。
  \item 合成データおよびCASASスマートホームセンサデータを用いた4つの実験により、実用的な処理速度と
        解釈可能な出力を実証する。
\end{itemize}
